\section{Experimental Setup}
\label{sec:setup}

\subsection{Simulated platforms}

The simulated platforms reproduce behaviours that real platforms are reported to have: the ladder is
selected by the upload resolution, ingest denoising and capped-CRF rate control depend on the tier,
and the viewer receives a fixed 720p rendition (\cref{tab:platforms}). The ladder rules are
plausible but invented; the tier gain is therefore real in the simulator by construction, and its
size says nothing about a real platform. The optimiser only calls
$\text{process}(\text{upload})\to\text{playback}$.

\begin{table}[t]
  \centering
  \caption{Hidden tiers of the simulated platforms. Rate control is capped CRF (maximum rate and
  buffer equal to the cap); playback is 720p.}
  \label{tab:platforms}
  \small
  \begin{tabular}{@{}lllrrl@{}}
    \toprule
    Platform & Upload height & Codec & CRF & Cap (kb/s) & Ingest denoise \\
    \midrule
    \multirow{3}{*}{\texttt{ugc\_strict}}
      & $\ge 1440$ & x265 (fast)   & 26 & 2000 & 1.5:1.5:4:4 \\
      & $\ge 1080$ & x264 (medium) & 25 & 1600 & 2:2:6:6 \\
      & otherwise  & x264 (medium) & 27 & 1200 & 3:3:8:8 \\
    \midrule
    \multirow{2}{*}{\texttt{ugc\_av1}}
      & $\ge 1440$ & SVT-AV1 (preset 10) & 36 & 1800 & none \\
      & otherwise  & x264 (medium)       & 24 & 1800 & 1:1:3:3 \\
    \bottomrule
  \end{tabular}
\end{table}

\subsection{Clips}

All clips are normalised to the same master format: $1280\times720$, 30\,fps, 3\,s (90 frames),
losslessly encoded.

\paragraph{P0 set.}
Four natural sequences from the SVT multi-format test set~\cite{haglund2006svt}, distributed by
Xiph.org (park\_joy, old\_town\_cross, ducks\_take\_off, in\_to\_tree), and three synthetic clips
(a slowly moving gradient, a zooming fractal and a test pattern with text). Searches train on
park\_joy and old\_town\_cross; the other clips are held out.

\paragraph{P1 set.}
77 segments from 47 sources, each with a recorded license: 53 natural segments (Xiph.org: SVT, TUM,
NTIA/ITS and public-domain video-conference sequences), 14 animation segments (Blender open movies
and rendered e-card sequences), 4 rendered screen-content segments and 6 synthetic segments.
Datasets without clear license terms were excluded. The per-source terms are listed with the code;
several sources are for research use only. The figures in this paper are charts and show no frames
of the test sequences. \todo{cite the sources individually (SVT, TUM, NTIA/ITS, Blender Foundation,
HDgreetings, Xiph.org video-conference sequences).}
Frame rates of 50/60\,fps are reduced to 30\,fps by dropping frames; 24/25\,fps sources are
retimed frame for frame, which speeds their motion up by 1.2--1.25$\times$.

\paragraph{Content classes.}
For each clip we compute the ITU-T P.910 spatial and temporal information (SI, TI; maximum over
frames)~\cite{itup910}, the fraction of $16\times16$ luma blocks with standard deviation below 2
(\emph{flat fraction}) and the fraction of luma values below 40 (\emph{dark fraction}). Classes are
assigned by a fixed rule, in this order:
\begin{itemize}
  \item \textbf{D} (graphics, text, screen, animation): source type is not natural;
  \item \textbf{C} (flat or dark): flat fraction $\ge 0.5$ or dark fraction $\ge 0.5$;
  \item \textbf{B} (dense texture or high motion): SI or TI at or above its 66th percentile over the
        natural clips not in C (SI $\ge 109.3$ or TI $\ge 33.1$);
  \item \textbf{A} (natural, low motion, moderate detail): otherwise.
\end{itemize}
Splits are 50/25/25\,\% (train, development, confirmation) per class, grouped by source across
classes. \Cref{tab:classes} gives the counts. Class C received a single clip; it is reported but
excluded from training and from every statistical test, and its gate entry is fixed in advance to
the tier-only configuration. The rule, thresholds, splits and gate entry for C were committed to the
repository before any training.

\begin{table}[t]
  \centering
  \caption{Content classes and splits of the P1 set (frozen before training).}
  \label{tab:classes}
  \begin{tabular}{@{}lrrrrr@{}}
    \toprule
    Class & Train & Dev & Confirm & Report only & Sources \\
    \midrule
    A: natural, low motion, moderate detail & 13 & 7 & 7 & -- & 17 \\
    B: dense texture or high motion         & 13 & 6 & 6 & -- & 14 \\
    C: flat or dark                         & -- & -- & -- & 1 & 1 \\
    D: graphics, text, screen, animation    & 12 & 6 & 6 & -- & 17 \\
    \bottomrule
  \end{tabular}
\end{table}

\subsection{Hypotheses (P1, pre-registered)}

\begin{itemize}
  \item \textbf{H1} (confirmatory): on class A, a learned pre-compensation beats the class-tuned
        sharpening-only configuration (\cref{sec:decision}).
  \item \textbf{H2} (exploratory): on class B it does not, because the emulator is least accurate on
        dense texture.
  \item \textbf{H3} (exploratory): on classes C and D every filter, the learned one included, is worse
        than tier-only, so the gate should disable filtering. H3 is not testable for C ($n=1$).
\end{itemize}

\paragraph{Baselines.}
Naive upload; tier-only ($\times2$); the global sharpening-only configuration found in P0; and a
class-tuned sharpening-only configuration (GP-BO, 15 trials on the class's training clips, with every
knob except the unsharp amount and radius frozen). The class-tuned sharpening-only configuration is
the baseline a learned pre-compensation must beat.

\subsection{Follow-up experiments}
Each follow-up was pre-registered in the repository (plan committed before the run) and uses the
frozen P1 classes and splits, judges and bootstrap intervals.
\begin{itemize}
  \item \textbf{Anchored estimator (E7).} The P1 global $G$ retrained with the estimator of
        \cref{eq:anchored}, with the emulator of the P1 run of the same seed frozen and every other
        setting unchanged; six anchors (every 250 steps) and, as a low-budget variant, one anchor at
        the identity; three seeds each.
  \item \textbf{$G$ on top of $\theta_A$ (E9).} $\theta_A$ placed in front of the black box, a new
        emulator (156 pairs) and $G$ with six anchors; three seeds.
  \item \textbf{Fixed recipes (E11)} (\cref{sec:recipes}) on all 77 clips.
  \item \textbf{Transfer (E14)}: the sharpening-only configurations and P0's full configurations
        evaluated on the other simulated platform.
  \item \textbf{Validation size (E10)}: subsampling of the E11 results on both platforms, 10{,}000
        clip subsets per size.
  \item \textbf{Passthrough rule (E12)}: \texttt{ugc\_strict} plus an invented rule that serves
        H.264 uploads of at most 720p and 2000\,kb/s (actual average) unchanged; GP-BO over upscale,
        unsharp amount and upload bitrate (300--80{,}000\,kb/s, log scale), 30 trials, two seeds.
\end{itemize}

\subsection{Compute}
One workstation with a 6-core CPU and an RTX 4060 GPU. A trial costs about 15--40\,s per clip.
\todo{report total GPU and CPU hours per experiment.}
